\chapter{CONCLUSIONS AND DISCUSSION}

\section{Summary}
This project designed and assessed an HSI classification system that allows direct comparison of supervised, semi-supervised and federated learning.

The comparison system used the same preprocessing pipeline, three benchmark datasets (Indian Pines, Pavia University, and Salinas) and three model architectures (2D-CNN, 3D-CNN, and Vision Transformer).

Spectral data that exceeds 200 bands (and is raw) is submitted to a series of treatments that include global min-max normalization, Gaussian spatial filtering, PCA that reduces the filtration to 30 components, and segmentation into 15 × 15 × 30 patches. From these, class proportions are maintained in a stratified 90/10 train-test split. In the semi-supervised model, self-training is combined with a curriculum that reduces the self-training model’s error propagation by introducing pseudo-labels that are least confident and/or least certain, in a progressive manner. The federated model mimics a three-client system in which each of the datasets is a local data silo. In this model, the FedAvg protocol is used to aggregate local copies of models in a central system without any transfer of original data.

\section{Key Findings}

\subsection{Semi-Supervised Learning Effectiveness}
Semi-supervised learning is more successfully applied to larger datasets.
\begin{enumerate}
    \item For Pavia University and Salinas, which contain large unlabeled data pools, semi-supervised models received competitive accuracy (86--97\%).
    \item For the smaller, more class-imbalanced Indian Pines set, accuracies were around 54--74\%, indicating the confidence threshold and the count of pseudo-label iterations require more tuning to the specific dataset.
    \item In all studies, 2D-CNN has the best semi-supervised average at 89.22\%, suggesting that simpler convolutional architectures are more forgiving of pseudo-label noise than deeper or attention-based ones.
    \item Curriculum learning does not remove the overfitting from noisy pseudo-labels, but does help, especially when the first labeled set is sufficiently large to create good pseudo-labels.
\end{enumerate}

\subsection{Federated Learning Privacy-Performance Trade-off}
Federated learning provides near-supervised accuracy, but introduces a communication burden.
\begin{enumerate}
    \item Federated 3D-CNN has a 97.36\% average overall accuracy, which only represents a $-2.55$ percentage point loss against its centralized supervised model (99.91\%).
    \item On Pavia University and Salinas, federated models achieve 99--100\% accuracy, which shows that FedAvg is applicable when local data is somewhat heterogeneous.
    \item The Indian Pines gap is larger (85--96\% federated vs.\ 99--100\% supervised), and shows the difficulty of getting meaningful weight updates from a small local dataset after 50 communication rounds.
    \item For organizations that cannot centralize raw spectral data, federated learning is a valuable option, since the accuracy cost is small and the privacy is protected.
\end{enumerate}

\subsection{Architecture-Specific Insights}
\begin{itemize}
    \item \textbf{2D-CNN:} Leads overall accuracy in every paradigm (supervised: 99.96\%, semi-supervised: 89.22\%, federated: 98.72\%). After PCA reduces the spectral axis to 30 components, 2D spatial convolutions are enough. Volumetric operations only add parameters with no proportional return in accuracy.

    \item \textbf{3D-CNN:} Placed second with consistency (supervised: 99.91\%, semi-supervised: 86.92\%, federated: 97.36\%). Modeling inter-band correlation with volumetric kernels results in better handling of noisy pseudo-labels and federated weight averaging.

    \item \textbf{Vision Transformer:} Competitive at full supervision (99.72\%) but most sensitive to distribution shift: average semi-supervised drops at 77.41\% and at 94.11\% federated. Global self-attention, although strong at full supervision with clean data, becomes easily destabilized with pseudo-labels noise and data distribution heterogeneity on clients.

    \item \textbf{Cross-paradigm pattern:} All architectures show significant deterioration to semi-supervised relative to both supervised and federated baselines. This suggests that in semi-supervised learning, the quality of pseudo-labels constitute the main bottleneck, as opposed to the architectures’ capacity.
\end{itemize}

\section{Project Objectives and Achievement Level}

\subsection{Original Objectives}
The project aimed to address the following goals.
\begin{enumerate}
    \item Create normalization and dimensionality reduction pipelines for high-dimensional hyperspectral data.
    \item Establish baseline supervised learning systems based upon different deep learning architectures (CNN, ViT).
    \item Integrate semi-supervised learning systems based on self-training and curriculum learning to overcome the problem of scarce labeling.
    \item Introduce federated learning to enable privacy concerns to be addressed during collaborative learning across distributed data.
    \item Create a reactive GUI with multithreaded inference.
    \item Provide benchmark datasets with an overall accuracy exceeding 95%.
\end{enumerate}

\subsection{Achievement Level}
All of the aforementioned goals were achieved (or more than were).

\begin{table}[H]
    \centering
    \caption{Project Objectives and Achievement Level}
    \label{tab:objectives}
    \renewcommand{\arraystretch}{1.3}
    \resizebox{\textwidth}{!}{%
    \begin{tabular}{|l|c|c|}
        \hline
        \textbf{Objective} & \textbf{Target} & \textbf{Achieved} \\ \hline
        Preprocessing Pipeline & Normalization + Filtering + PCA & $\checkmark$ Fully Implemented \\ \hline
        Supervised Baselines & Multiple architectures & $\checkmark$ 2D-CNN, 3D-CNN, ViT \\ \hline
        Semi-Supervised Learning & Self-training + CL & $\checkmark$ Implemented (54--97\% range) \\ \hline
        Federated Learning & Multi-client aggregation & $\checkmark$ FedAvg Protocol (85--99\% range) \\ \hline
        GUI with Multithreading & Python (PyQt6) + Multithreading & $\checkmark$ Fully Operational \\ \hline
        Target Accuracy ($>$ 95\%) & Supervised requirement & $\checkmark$ Achieved (99.72--99.96\%) \\ \hline
    \end{tabular}}
\end{table}

\section{Strengths}

\subsection{Methodological Strengths}
\begin{itemize}
    \item \textbf{Unified comparative framework:} The ability to evaluate three paradigms under the same preprocessing, splitting, and metrics is uncommon in the literature related to HSI; results are truly comparable.

    \item \textbf{Curriculum Learning:} Combining self-training and the presentation of samples in a scaffolded format significantly alleviates pseudo-label overfitting and the task suddenly ceases to converge to deteriorated local minima.

    \item \textbf{Privacy-preservation federated framework:} The federated framework demonstrates that pooling sensitive spectral data is not a prerequisite for competitive classification results. This is directly related to remote sensing projects involving multiple institutions.

    \item \textbf{Generalizable preprocessing:} The three-step pipeline (normalization, Gaussian filtering, fixed-component PCA) has been validated with three datasets with different class counts, spatial resolution, and bands. The three steps of this data-driven pipeline are generalizable.
\end{itemize}

\subsection{Implementation Strengths}
\begin{itemize}
    \item \textbf{Three-architecture support:} The flexibility of the surrounding pipeline means users may choose between a 2D-CNN, 3D-CNN, and ViT based on preferences for accuracy and efficiency.

    \item \textbf{GPU-accelerated training:} Large datasets (N > 50,000) can be trained using (+) Based on the benchmarking data presented in the paper, the framework achieves very high training efficiency with the use of Google Colab GPU.

    \item \textbf{Responsive multithreaded GUI:} PyQt6 decouples inference from the UI event loop and allows the GUI to remain active during long lasting classification.

    \item \textbf{Reproducibility:} Fixed randomization and a stratified split guarantee reproducible results from an experimental run and do not require further (manual) tuning.
\end{itemize}

\section{Limitations and Weaknesses}

\subsection{Methodological Limitations}
\begin{itemize}
    \item \textbf{Simulated federated scenario:} Having three local datasets on one machine simulates three clients. True deployments have unbalanced datasets, considerable latency, and stragglers, therefore, the accuracy gap would be expected to be larger.

    \item \textbf{Fixed PCA components:} The empirical decision to fix 30 PCA components could be evolved to a more sophisticated choice of variance or information retention for the given scene.

    \item \textbf{Confidence threshold rigidity:} All datasets used the same threshold for self-training (confidence > 0.95). Sensitivity tests on class-imbalanced datasets could demonstrate the same degree of flexibility in the semi-supervised pipeline.

    \item \textbf{Single curriculum schedule:} Difficulty-only curriculum ordering has been used exclusively. Reverse order, flexible pace, and other (class size driven) uncertainty curricula are untested and may have significant effects based on class size.
\end{itemize}

\subsection{Experimental Limitations}
\begin{itemize}
    \item \textbf{Benchmark datasets only:} Generalizing to actual operational hyperspectral surveys with real acquisition noise and non-standard class taxonomies is impossible.

    \item \textbf{No concept drift analysis:} The distributed system assumes the data remains stationary. Remote sensing data has a temporal drift and would require concept drift in a production setting.

    \item \textbf{Communication cost not quantified:} While federated accuracy is presented, no measurements of bandwidth usage or one-way latency is presented, which are both crucial for assessing the practicality of the system.

    \item \textbf{Adversarial robustness not assessed:} Operational deployment requires the analysis of the perturbed distributions of sensors and a sample data set; both are lacked.
\end{itemize}

\subsection{Implementation Limitations}
\begin{itemize}
    \item \textbf{No model compression:} Weight pruning and quantization are absent.

    \item \textbf{Single-GPU training:} Data parallelism is absent; training is limited to a single GPU.

    \item \textbf{Restricted hyperparameter tuning:} The approach is a coarse grid search, which provides little estimate improvement. Confidence thresholds in semi-supervised learning may benefit significantly from hyperparameter tuning.
\end{itemize}

\section{Future Work}

\subsection{Methodological Extensions}
\begin{enumerate}
    \item \textbf{Federated continual learning:} Imparts the distributed system the capability to adapt to the temporal distribution shifts of hyperspectral data without requiring a complete system refresh.

    \item \textbf{Heterogeneous federated aggregation:} Introduces hierarchical aggregations (weighted aggregation) to accommodate differing data distributions of the clients.

    \item \textbf{Adaptive difficulty:} Rather than the fixed ordering of training sample difficulties, allow the model to determine the sample order based on its current state.

    \item \textbf{Few-shot HSI:} Develop the ability to rapidly transfer learning in the context of few-shot HSI across different sensors and regions.
\end{enumerate}

\subsection{Architectural Enhancements}
\begin{enumerate}
    \item \textbf{Hybrid 2D-3D models:} Integrate 2D and 3D convolutions to find optimal trade-off between spectral-spatial expressiveness and number of model parameters; the promising outcomes of 2D-CNNs indicate that this may define the efficient frontier.

    \item \textbf{Graph Neural Networks:} Model spectral-spatial relationships as graphs to capture irregular spatial relationships which may be lost with traditional convolutional approaches.

    \item \textbf{Self-supervised pretraining:} Leverage contrastive or masked-reconstruction self-supervised pretraining on unlabeled hyperspectral data to build foundational representations which can then be improved through semi-supervised learning.

    \item \textbf{Multi-modal fusion:} Extend the workflow to integrate hyperspectral data with panchromatic images and laser data within those contexts where one modality falls short.
\end{enumerate}

\subsection{Experimental Directions}
\begin{enumerate}
    \item \textbf{Real-world federated deployment:} Test the federated learning framework with actual datasets from known multiple data sources/workplaces to validate that the simulated results will hold in real-world use cases (e.g., ESA, NASA).

    \item \textbf{Cross-sensor generalization:} Evaluate the federated learning framework with multiple disparate sensors (AVIRIS, Sentinel-2, DESIS) for an understanding of the extent of domain generalization and adaptation.

    \item \textbf{Adversarial robustness:} Perform comprehensive adversarial evaluations and create sensor noise defense frameworks for fully deployed classifiers.

    \item \textbf{Uncertainty quantification:} Implement Bayesian deep learning for hypothetical low integration risk and high-value remote sensing applications.
\end{enumerate}

\subsection{Practical Applications}
\begin{enumerate}
    \item \textbf{Edge deployment:} Use the combination of model pruning and quantization to achieve the goal of real-time data classification from nanosatellites.

    \item \textbf{Active learning interface:} Implement an active learning module in the interface to autonomously flag the most informative unlabeled pixels.

    \item \textbf{Interpretability tooling:} Integrate explainable AI (attenuation mapping, LIME, SHAP) to facilitate avenues for domain experts to perform thorough evaluations of model outputs prior to engagement.

    \item \textbf{GUI enhancements:} Annotation tools, quality assurance workflows, and real-time performance dashboards will be added to the interface based on the needs of the operational remote sensing teams.
\end{enumerate}

\section{Conclusion}
There is no ambiguity in the data. First, in federated learning, a small average reduction in accuracy (-3.1\%) offers a cost-benefit tradeoff for significant data privacy, especially in cases with large data. Second, in the case of semi-supervised learning, self-training, and curriculum learning, rapid scaling is achievable with competition on the Pavia University and Salinas datasets, but on the small, imbalanced Indian Pines set, dominated by noisy pseudo-labels, the system can’t keep pace – a loss of efficacy is expected. 2D-CNNs are the clear winners of all paradigms. One interesting thing to note is that, after PCA on the spectral data, the spatial data alone is sufficient for near-perfect classification. The Vision Transformer, though achieving the same level of accuracy as CNNs in full supervision, is the least resilient to the distribution shifts resulting from pseudo-labeling and federated weight averaging. What has been developed is a comprehensive system that integrates three paradigms, three architectures, and three datasets, all within a user-friendly multi-threaded GUI. Many examples of research can be found online, but what this research incorporates is an application that can be readily implemented. The focus of subsequent research needs to be on actual federated implementations and on curriculum techniques that are dynamically responsive. These techniques should be tailored to address the class imbalance in the datasets, as these are the most influential factors to the system’s performance.